% BAPC 2026 team reference document. Build with ./build.sh, test the code with tests/run.sh.
% TRD rules: max 25 single-sided A4 pages, university name top left, page number top right,
% readable from half a metre. Handwritten notes are allowed on the front of the pages.
\documentclass[9pt]{extarticle}
\usepackage[a4paper, landscape, left=7mm, right=7mm, top=14mm, bottom=7mm,
            headsep=3mm, headheight=12pt]{geometry}
\usepackage[T1]{fontenc}
\usepackage[utf8]{inputenc}
\usepackage[scaled=0.92]{helvet}
\renewcommand{\familydefault}{\sfdefault}
\usepackage[varqu, varl]{inconsolata}
\usepackage{amsmath, amssymb}
\usepackage{xcolor, listings, multicol, fancyhdr, enumitem, booktabs, tabularx, array}
\usepackage[explicit]{titlesec}
\usepackage{titletoc}
\usepackage[most]{tcolorbox}
\usepackage{needspace}
\usepackage{microtype}
% in \texttt, << and >> would otherwise print as French quotes
\DisableLigatures[<,>]{encoding = T1, family = tt*}

% ---- Fill these in before printing ----
\newcommand{\university}{University of Twente}
\newcommand{\teamname}{Nth time's the charm}

% ---- Page header: university left, section in the middle, page number right ----
\pagestyle{fancy}
\fancyhf{}
\fancyhead[L]{\textbf{\university}\enspace|\enspace\teamname}
% the chapters on this page, e.g. "5 Dynamic programming -- 6 Graphs"
\makeatletter
\newcommand{\pagesections}{%
  \edef\nb@top{\expandafter\@leftmark\topmark\@empty\@empty}%
  \edef\nb@last{\leftmark}%
  \ifx\nb@top\@empty \nb@last
  \else\ifx\nb@top\nb@last \nb@last
  \else \nb@top\ \textendash\ \nb@last \fi\fi}
\makeatother
\fancyhead[C]{\ifnum\value{page}=1 Contents\else\pagesections\fi}
\fancyhead[R]{\textbf{\large\thepage}}
\renewcommand{\headrulewidth}{0.4pt}
\fancypagestyle{notes}{\fancyhf{}%
  \fancyhead[L]{\textbf{\university}\enspace|\enspace\teamname}%
  \fancyhead[C]{Notes}\fancyhead[R]{\textbf{\large\thepage}}}

\setlength{\columnsep}{4mm}
\setlength{\columnseprule}{0.3pt}
\setlength{\parindent}{0pt}
\setlength{\parskip}{1.5pt}
\setlist{nosep, leftmargin=1em, label=\textbullet}

% ---- Headings ----
\titleformat{\section}{\normalfont\Large\bfseries}{}{0pt}{\thesection\quad #1}[{\titlerule[1pt]}]
\titleformat{name=\section, numberless}{\normalfont\Large\bfseries}{}{0pt}{#1}[{\titlerule[1pt]}]
\titlespacing*{\section}{0pt}{6pt}{3pt}
\titleformat{\subsection}{\normalfont\bfseries}{}{0pt}%
  {\colorbox{black!15}{\parbox{\dimexpr\linewidth-2\fboxsep}{\strut\thesubsection\quad #1}}}
\titlespacing*{\subsection}{0pt}{5pt}{2pt}
\renewcommand{\sectionmark}[1]{\markboth{\thesection\enspace #1}{}}

% ---- Table of contents ----
\titlecontents{section}[0pt]{\addvspace{3pt}\bfseries}{\thecontentslabel\enspace}{}{\hfill\contentspage}
\newcommand{\tocend}{\par}  % \par is not allowed directly in the argument
\titlecontents*{subsection}[0pt]{\small\raggedright}{}{}{~\thecontentspage}[\enspace\textbullet\enspace][\tocend]

% ---- Code ----
\lstdefinestyle{cpp}{
  language=C++, morekeywords={ll, vi, pii, rep, all, auto},
  basicstyle=\ttfamily\fontsize{7.4}{8.3}\selectfont,
  keywordstyle=\bfseries,
  commentstyle=\color{black!60},
  stringstyle=\color{black!80},
  columns=fullflexible, keepspaces=true,
  breaklines=true, postbreak=\mbox{$\hookrightarrow$\space},
  tabsize=4, showstringspaces=false, upquote=true,
  frame=leftline, framerule=1pt, rulecolor=\color{black!40}, framesep=3pt, xleftmargin=4pt,
  aboveskip=2pt, belowskip=3pt,
}
% \code{name} prints code/cpp/name.cpp (keep its lines at most 68 characters)
% every snippet is one unbreakable block (the longest, 55 lines, fits a column)
\newcommand{\keep}[1]{\par\noindent\begin{minipage}{\linewidth}#1\end{minipage}\par}
\newcommand{\code}[1]{\keep{\lstinputlisting[style=cpp]{code/cpp/#1.cpp}}}
% \pycode{name} and \shcode{name}: code/py/name.py and code/sh/name.sh
\newcommand{\pycode}[1]{\keep{\lstinputlisting[style=cpp, language=Python]{code/py/#1.py}}}
\newcommand{\shcode}[1]{\keep{\lstinputlisting[style=cpp, language=bash]{code/sh/#1.sh}}}
% a shell command, e.g. how to compile
\newcommand{\shell}[1]{\par{\raggedright\small\ttfamily #1\par}}

% ---- "Which technique?" lines: clue, what to try, page ----
\newcommand{\techrow}[3]{\par\noindent
  \parbox[t]{0.50\linewidth}{\raggedright #1}\hfill
  \parbox[t]{0.40\linewidth}{\raggedright #2}\hfill
  \makebox[1.4em][r]{#3}\par\vspace{0.6pt}}
\newcommand{\techhead}{\techrow{\textbf{the statement says}}{\textbf{try}}{\textbf{p.}}%
  \vspace{-1pt}\noindent\rule{\linewidth}{0.4pt}\par}
\newcommand{\tech}[3]{\techrow{#1}{#2}{\pageref{t:#3}}}

% ---- Topics ----
% \topic{Dijkstra}{$O(m \log n)$}: a heading with its running time on the right
\newcommand{\topic}[2]{\subsection[#1]{#1\hfill #2}}
% \use{...}: one line saying when to use it, in a grey box
\newtcolorbox{tip}{enhanced, colback=black!5, colframe=black!5, boxrule=0pt, sharp corners,
  left=3pt, right=3pt, top=1pt, bottom=1pt, before skip=1pt, after skip=1pt,
  before upper=\raggedright, after={\par\nobreak},
  borderline west={1.5pt}{0pt}{black!60}}
\newcommand{\use}[1]{\begin{tip}#1\end{tip}}

\begin{document}
\begin{multicols*}{3}
\raggedcolumns
\raggedright   % narrow columns: no stretched spaces
\ifdefined\onlysection
  \input{sections/\onlysection}   % ./preview.sh name: one section only
\else
  \tableofcontents
  \section{Contest}

\subsection{How big can $n$ be?}
About $10^8$ simple steps per second.\par
\begin{tabularx}{\linewidth}{@{}rlX@{}}
\toprule
$n \le$ & time & typical idea\\
\midrule
10 & $O(n!)$ & every order (\texttt{next\_permutation})\\
20 & $O(2^n \cdot n)$ & every subset, bitmask DP\\
500 & $O(n^3)$ & three loops, Floyd--Warshall\\
5\,000 & $O(n^2)$ & two loops, $n \times n$ DP\\
$10^6$ & $O(n \log n)$ & sort, binary search, Dijkstra, set/map\\
$10^8$ & $O(n)$ & one loop, prefix sums\\
$10^{18}$ & $O(\log n)$ & formula, binary search on the answer\\
\bottomrule
\end{tabularx}\par
\texttt{int} $\approx 2.1\cdot10^9$, \texttt{long long} $\approx 9.2\cdot10^{18}$,
\texttt{double} $\approx 15$ digits.

\subsection{Which technique?}
% one line per clue, so the list can continue in the next column
\techhead
\tech{fewest moves / steps, maze}{BFS, grid BFS}{bfs}
\tech{cheapest route, weighted roads}{Dijkstra (shortest path table)}{sp}
\tech{``smallest $x$ such that it works''}{binary search on the answer}{binsearch}
\tech{best total over choices, $n \le 5000$}{DP}{dp}
\tech{number of ways (mod $10^9+7$)}{DP, counting formulas}{counting}
\tech{choose items, weight limit}{knapsack}{knapsack}
\tech{$n \le 20$}{all subsets, bitmask DP}{subsets}
\tech{$n \le 10$, find the best order}{all orders}{perms}
\tech{intervals, meetings, schedules}{sort by end, greedy}{intervals}
\tech{rooms needed at the same time}{sweep $+1 / -1$}{sweep}
\tech{sort and take, pair up}{greedy moves}{greedy}
\tech{``connected?'', groups merge}{Union-Find}{dsu}
\tech{``$a$ before $b$'', dependencies}{topological sort}{topo}
\tech{connect everything cheaply}{Kruskal}{mst}
\tech{tree, distances between nodes}{tree facts, LCA}{trees}
\tech{directed cycles, ``reach each other''}{SCC}{scc}
\tech{same step repeated $10^{18}$ times}{binary lifting, cycles}{lca}
\tech{sum of a range, many queries}{prefix sums}{prefix}
\tech{add to a range, many times}{difference array}{diff}
\tech{best window / subarray}{two pointers}{twoptr}
\tech{nearest bigger to the left / right}{monotonic stack}{mono}
\tech{up then down, few questions}{peak search}{peak}
\tech{``interactive''}{flush, binary search}{interactive}
\tech{``modulo $10^9 + 7$''}{modular arithmetic}{mod}
\tech{repeats every $a$ and every $b$}{gcd / lcm, CRT}{gcd}
\tech{primes, divisors, factors}{sieve, factorising}{sieve}
\tech{recurrence, $n$ up to $10^{18}$}{matrix power}{matrix}
\tech{digits of numbers up to $10^{18}$}{digit DP}{digit}
\tech{hundreds of digits}{big numbers, Python}{big}
\tech{probability, expected value}{expected-value DP}{expect}
\tech{two players, optimal play}{games}{games}
\tech{substrings, patterns, words}{string tool table}{strings}
\tech{points, lines, polygons, angles}{geometry rules}{geo}
\tech{clocks, dates}{time and dates}{time}
\tech{rotate / mirror a grid}{grids}{grid}
\tech{text with commas, odd input}{splitting text}{parse}

\subsection{Template}
\code{template}

\subsection{Terminal setup}
\shcode{contest}

\subsection{Reading input}
\code{io}

  \section{C++ quick reference}

\subsection{Sorting}
\code{stl_sort}

\subsection{Useful algorithms}
\code{stl_algo}

\subsection{Containers}
\use{Sorted / next bigger: \texttt{set}, \texttt{map} ($\log n$). Lookup only:
\texttt{unordered\_map}. Always the min or max: \texttt{priority\_queue}. Both ends:
\texttt{deque}.}
\code{stl_containers}

\subsection{Strings and characters}
\code{stl_strings}

\subsection{Bits}
\code{stl_bits}

\subsection{Lambdas}
\code{stl_lambda}

\topic{Ordered set}{$O(\log n)$}
\use{The $k$-th smallest, or how many are below $x$, while inserting and erasing.}
\code{stl_oset}

\topic{Hash for pair keys}{$O(1)$}
\use{\texttt{unordered\_map} with \texttt{pair} keys does not compile without it.}
\code{stl_pairhash}

  \section{Simulation and parsing}

\topic{Grids}{$O(R \cdot C)$}\label{t:grid}
\use{Rotate, transpose, mirror; 8 neighbours; wrap around the edges.}
\code{adhoc_grid}

\topic{Time and dates}{$O(1)$}\label{t:time}
\use{Clock times as minutes, leap years, days between dates, day of the week.}
\code{adhoc_time}

\topic{Splitting text}{$O(n)$}\label{t:parse}
\use{Fields separated by commas or other characters; numbers hidden in text.}
\code{adhoc_parse}

\topic{Number bases}{$O(\log x)$}
\use{Binary, hex, base 36 and back.}
\code{adhoc_base}

  \section{Sorting and greedy}

\subsection{Greedy moves}\label{t:greedy}
\begin{tip}
\begin{itemize}
\item Sort, then take the biggest (or smallest) first.
\item Pair the smallest with the largest.
\item End every piece as early as possible.
\item Give each person the smallest thing that is big enough.
\item Not sure? Break it with 2--3 tiny cases or a brute force. Fails and $n$ small: DP.
\end{itemize}
\end{tip}

\topic{Finding the sort order}{$O(n \log n)$}
\use{Compare the cost of ``$a$ then $b$'' with ``$b$ then $a$''. Below: $a$ first iff
$t_a w_b < t_b w_a$. Biggest glued number: sort strings by \texttt{a + b > b + a}.}
\code{greedy_exchange}

\topic{Most non-overlapping intervals}{$O(n \log n)$}\label{t:intervals}
\use{Take the interval that ends first. Fewest points hitting every closed $[s, e]$: the same
code with end $e + 1$.}
\code{greedy_intervals}

\topic{Merging intervals}{$O(n \log n)$}
\use{Union of intervals: length covered, gaps, ``is $x$ covered?''.}
\code{greedy_merge}

\topic{Most intervals at the same time}{$O(n \log n)$}\label{t:sweep}
\use{Rooms or machines needed at once: $+1$ at a start, $-1$ at an end, sort, running count.}
\code{greedy_sweep}

\topic{Jobs with deadlines}{$O(n \log n)$}
\use{Most profit or most jobs in time: by deadline, take every job; when they no longer fit,
drop the worst one taken.}
\code{greedy_deadlines}

\topic{Matching with the smallest that fits}{$O((n + m) \log m)$}
\use{Person $i$ accepts sizes in $[lo_i, hi_i]$, most pairs: smallest $hi$ first, give it the
smallest item $\ge lo$. Only ``at least $lo$'': $hi = $ \texttt{INF}.}
\code{greedy_matching}

\topic{The answer is one of the input values}{$O(n \log n)$}
\use{Choose one number (price, limit, time): try every input value (and $\pm 1$, 0).
Smallest $\sum |x - a_i|$: $x$ is the median.}
\code{greedy_candidates}

  \section{Searching and arrays}

\topic{Binary search on the answer}{$O(\log(\mathit{hi} - \mathit{lo}))$ checks}\label{t:binsearch}
\use{``Smallest $t$ such that \dots'' and checking one $t$ is easy: works for $t$ $\Rightarrow$
works for $t + 1$. Write \texttt{ok(t)}.}
\code{binary_search}

\topic{Binary search on real numbers}{100 steps}
\use{The answer is a real number and \texttt{ok} is monotone.}
\code{binary_search_real}

\topic{Finding a peak}{$O(\log n)$}\label{t:peak}
\use{Goes up then down (or down then up), find the top. Often interactive. Integers: binary
search on ``is the next value lower?''; reals: ternary search.}
\code{peak_int}
\code{ternary}

\subsection{Searching in sorted data}
\code{lower_bound}

\topic{Prefix sums}{$O(1)$ per range}\label{t:prefix}
\use{Sum of many ranges; 2D: sum of any rectangle.}
\code{prefix}
\code{prefix2d}

\topic{Difference array}{$O(1)$ per update}\label{t:diff}
\use{Many ``add $v$ to $a[l..r]$'', the final values only at the end.}
\code{diff_array}

\topic{Two pointers / sliding window}{$O(n)$}\label{t:twoptr}
\use{Best subarray with a rule that stays true when the window shrinks (sum $\le K$ with
values $\ge 0$, at most $k$ different letters).}
\code{two_pointers}

\topic{Coordinate compression}{$O(n \log n)$}
\use{Huge values ($10^9$), few of them: replace by ranks to use as indices.}
\code{compress}

\topic{Monotonic stack}{$O(n)$}\label{t:mono}
\use{Nearest bigger element to the right. Left side: loop backwards. Smaller: flip the
comparison.}
\code{mono_stack}

\topic{Largest rectangle in a histogram}{$O(n)$}
\use{Largest free rectangle in a grid: per row, $h[c]$ = free cells straight above; run on
every row, $O(R \cdot C)$.}
\code{histogram}

\topic{Sliding window minimum}{$O(n)$}
\use{Min (or max) of every window of $k$, e.g.\ \texttt{dp[i] = min(dp[i-k..i-1]) + c}.}
\code{sliding_min}

  \section{Dynamic programming}

\subsection{How to build a DP}\label{t:dp}
\begin{tip}
\begin{itemize}
\item \textbf{State:} what you must know to finish the rest. \texttt{dp[i]}; \texttt{dp[i][j]}
  (two strings, a grid); \texttt{dp[l][r]} (a piece); add weight used, last item, speed\dots
\item \textbf{Transition:} try every last choice: \texttt{dp[i] = best(cost + dp[smaller])}.
\item \textbf{Base:} \texttt{dp[0] = 0}. Impossible: \texttt{INF} (min), \texttt{-INF} (max),
  0 (counting).
\item \textbf{Order:} fill a state after every state it reads, or memoize.
\item \textbf{Fast enough:} states $\times$ choices $\le 10^8$. Row $i$ only reads row
  $i - 1$: keep two rows.
\item \textbf{The choices:} store \texttt{from[state]}, walk back.
  \textbf{Counting:} \texttt{\% MOD} after every addition.
\item \textbf{Games:} winning iff some move leads to a losing position.
\end{itemize}
\end{tip}

\topic{Memoized recursion}{states $\times$ choices}
\use{The recursion is easy but solves the same subproblem again and again: store each answer.}
\code{dp_memo}

\topic{Best subarray sum}{$O(n)$}
\use{Largest sum of consecutive elements. Smallest: flip the signs.}
\code{dp_kadane}

\topic{0/1 knapsack and subset sum}{$O(n \cdot W)$}\label{t:knapsack}
\use{Each item at most once, weight $\le W$, most value. Subset sum: which totals can be made
(bitset: 64$\times$ faster). Two halves as equal as possible: the largest
$s \le \mathit{total}/2$ with \texttt{can[s]}.}
\code{dp_knapsack}
\code{dp_subset_sum}

\topic{Coin change}{$O(n \cdot S)$}
\use{Every coin as often as you like: fewest coins, number of ways.
Greedy fails: coins 1, 3, 4 for 6.}
\code{dp_coins}

\topic{Longest increasing subsequence}{$O(n \log n)$}
\use{Fewest deletions to make it sorted: $n - \mathrm{LIS}$. Boxes in boxes: sort by the
first value, ties by the second descending, then LIS on the second.}
\code{dp_lis}

\topic{Longest common subsequence}{$O(n \cdot m)$}
\use{The longest subsequence two strings (or arrays) share.}
\code{dp_lcs}

\topic{Edit distance}{$O(n \cdot m)$}
\use{Fewest insertions, deletions and replacements that turn $a$ into $b$.}
\code{dp_edit}

\topic{Paths in a grid}{$O(R \cdot C)$}
\use{Only moves right or down: count paths, most coins. Moves in every direction: BFS or
Dijkstra.}
\code{dp_grid}

\topic{Cutting a string into words}{$O(|s| \cdot L)$}
\use{Split a string into words from a list: possible? how many ways? fewest pieces
(\texttt{min} instead of \texttt{+})?}
\code{dp_wordbreak}

\topic{Cutting an array into $k$ parts}{$O(k \cdot n^2)$}
\use{Consecutive groups, each with a cost: cheapest total. Any number of groups: drop $j$.}
\code{dp_split}

\topic{Range sums inside a DP}{$O(n)$}
\use{\texttt{dp[i]} adds up a range of earlier states: keep a prefix sum of \texttt{dp}.}
\code{dp_prefix}

\topic{Weighted intervals}{$O(n \log n)$}
\use{Non-overlapping intervals with the most total value. A pass from day $d$ covering $L$
days: the next state is \texttt{lower\_bound}$(d + L)$.}
\code{dp_weighted}

\topic{Interval DP}{$O(n^3)$}
\use{Merging neighbours, cutting a stick, brackets, taking from both ends:
\texttt{dp[l][r]}, try every split point. $n \le 400$.}
\code{dp_interval}

\topic{Bitmask DP}{$O(2^n \cdot n^2)$}\label{t:bitmask}
\use{$n \le 20$ and you must remember which ones are used: visit all (TSP, $n \le 18$),
jobs to workers ($O(2^n \cdot n)$). Narrow grid: row by row, the state is a mask of the row.}
\code{dp_bitmask}
\code{dp_assign}

\topic{Digit DP}{$O(\mathit{digits} \cdot \mathit{states} \cdot 10)$}\label{t:digit}
\use{Count $x$ in $[A, B]$ with a digit property: \texttt{countUpTo(B) - countUpTo(A - 1)}.}
\code{dp_digits}

  \section{Graphs}

\subsection{Which shortest path?}\label{t:sp}
\begin{tabularx}{\linewidth}{@{}Xll@{}}
\toprule
the graph & use & time\\
\midrule
every edge costs 1 & BFS & $O(n+m)$\\
grid, every step costs 1 & grid BFS & $O(R \cdot C)$\\
costs only 0 or 1 & 0-1 BFS & $O(n+m)$\\
all costs $\ge 0$ & Dijkstra & $O(m \log n)$\\
some costs $< 0$ & Bellman--Ford & $O(n \cdot m)$\\
all pairs, $n \le 500$ & Floyd--Warshall & $O(n^3)$\\
no cycles (DAG) & DP in topological order & $O(n+m)$\\
\bottomrule
\end{tabularx}
\begin{tip}
\begin{itemize}
\item Many start points: push all of them with distance 0 (multi-source).
\item From every node to one target: reverse the edges, run from the target.
\item The path itself: keep \texttt{par[v]}, walk back from the end.
\item More than the position matters (direction, keys, turns, time mod $k$):
  the node is the state, e.g.\ \texttt{dist[r][c][dir]}.
\end{itemize}
\end{tip}

\subsection{Reading a graph}
\code{graph_types}
\code{graph_input}

\topic{BFS}{$O(n+m)$}\label{t:bfs}
\use{Fewest steps when every edge costs the same.}
\code{bfs}

\topic{Grid BFS}{$O(R \cdot C)$}\label{t:gridbfs}
\use{Maze or map of cells. 8 directions: add the diagonals to \texttt{dr}/\texttt{dc}.
Count islands: start a BFS from every free cell not seen yet.}
\code{grid_bfs}

\topic{0-1 BFS}{$O(n+m)$}
\use{Every cost is 0 or 1 (e.g.\ walking is free, breaking a wall costs 1).}
\code{bfs01}

\topic{Dijkstra}{$O(m \log n)$}
\use{Weighted edges, all weights $\ge 0$.}
\code{dijkstra}

\topic{Bellman--Ford}{$O(n \cdot m)$}
\use{Negative weights, or detect a negative cycle.}
\code{bellman_ford}

\topic{Floyd--Warshall}{$O(n^3)$}
\use{Distance between every pair, $n \le 500$. Negative edges are fine.}
\code{floyd}

\topic{DFS and connected components}{$O(n+m)$}
\use{Which nodes are connected, how many groups there are.}
\code{dfs}

\topic{Topological sort}{$O(n+m)$}\label{t:topo}
\use{``$a$ before $b$'' rules, cycle check, DP on a DAG: for $u$ in \texttt{order}, for
$u \to v$: \texttt{best[v] = max(best[v], best[u] + 1)}.}
\code{topo_sort}

\topic{Union-Find (DSU)}{$\approx O(1)$}\label{t:dsu}
\use{Merge groups and ask ``are $a$ and $b$ connected?'' while edges are added.}
\code{dsu}

\topic{Minimum spanning tree (Kruskal)}{$O(m \log m)$}\label{t:mst}
\use{Connect all nodes as cheaply as possible. A node may buy its own source for $c_v$:
add a node $X$ with edges $(X, v, c_v)$.}
\code{kruskal}

  \section{Trees and more graphs}

\subsection{Tree facts}\label{t:trees}
\begin{tip}
\begin{itemize}
\item $n$ nodes, $n - 1$ edges. Connected with $n - 1$ edges $\Rightarrow$ a tree.
\item Exactly one path between two nodes: one BFS gives all distances, also weighted.
\item Connected graph: $m - (n - 1)$ extra edges, each closes a cycle. Forest of $c$ trees:
  $n - c$ edges. Few extra edges: try all $2^k$ ways to handle them.
\item Edge above $v$ is used by $\mathit{sz}[v] \cdot (n - \mathit{sz}[v])$ paths. Summed over
  all edges: the sum of all pair distances.
\item Parent list with $p_i < i$: loop $i = n, \dots, 2$ for children before parents.
\item Trees are bipartite: colour by depth mod 2.
\end{itemize}
\end{tip}

\topic{Rooting a tree}{$O(n)$}
\use{Parent, depth, \texttt{order}. Backwards over \texttt{order}: children before parents
(sizes, tree DP). Forwards: parents first (rerooting). No recursion.}
\code{tree_root}
\code{tree_size}

\topic{DFS order}{$O(n)$}
\use{Enter / leave times. $u$ is in the subtree of $v$ iff
\texttt{tin[v] <= tin[u] < tout[v]}: a subtree is a range of an array.}
\code{tree_dfs}

\topic{Tree DP}{$O(n)$}
\use{The answer for $v$ follows from its children's answers. Fill backwards over
\texttt{order}; the answer is at the root.}
\code{tree_dp}

\topic{Diameter}{$O(n)$}
\use{Longest path: BFS from any node to the farthest $a$, BFS from $a$ to the farthest $b$.
The farthest node from any $v$ is $a$ or $b$.}
\code{tree_diameter}

\topic{LCA and binary lifting}{$O(\log n)$ per query}\label{t:lca}
\use{Distance between $a$ and $b$, the $k$-th node above $v$. Weighted:
$w[a] + w[b] - 2\,w[\mathrm{lca}]$. Every node has one next node (functional graph): where
after $k \le 10^{18}$ steps? \texttt{Jump}.}
\code{tree_lca}
\code{graph2_jump}

\topic{Rerooting}{$O(n)$}
\use{A tree DP answer for every node as the root. Moving the root from \texttt{par[v]} to
$v$: $\mathit{sz}[v]$ nodes get 1 closer, the other $n - \mathit{sz}[v]$ 1 farther.}
\code{tree_reroot}

\topic{Strongly connected components}{$O(n + m)$}\label{t:scc}
\use{Directed graph, nodes that reach each other become one node: the condensation is a DAG.
To reach everything you need a start in every component with in-degree 0.}
\code{graph2_scc}

\topic{Euler path}{$O(n + m)$}
\use{Every edge exactly once. Undirected: 0 or 2 odd-degree nodes (start at one).
Directed: in $=$ out everywhere, or one start (out $=$ in $+ 1$) and one end. Circuit: all
even / in $=$ out.}
\code{graph2_euler}

\topic{Bipartite check}{$O(n + m)$}
\use{Two groups, every edge between them: possible iff no odd cycle.}
\code{graph2_bipartite}

\topic{Bipartite matching}{$O(n \cdot m)$}
\use{Most pairs between two kinds of things. K\"onig: min vertex cover $=$ max matching;
max independent set $= n -$ max matching.}
\code{graph2_matching}

  \section{Math and number theory}

\subsection{Integer pitfalls}
\begin{tabularx}{\linewidth}{@{}Xll@{}}
\toprule
you want & not this & but this\\
\midrule
$\lceil a/b \rceil$ for $a, b > 0$ & \texttt{ceil(a/b)} & \texttt{(a+b-1)/b}\\
$\lfloor a/b \rfloor$, $a$ may be $< 0$ & \texttt{a/b} & \texttt{floorDiv(a,b)}\\
remainder in $[0, m)$ & \texttt{a\%m} & \texttt{md(a,m)}\\
$a \cdot b$ for two \texttt{int}s & \texttt{a*b} & \texttt{1LL*a*b}\\
a product up to $10^{36}$ & \texttt{ll} & \texttt{(\_\_int128)a*b}\\
$\lfloor \sqrt n \rfloor$ & \texttt{(ll)sqrt(n)} & \texttt{isqrt(n)}\\
$10^{18}$ or $3^k$ & \texttt{pow(10,18)} & a loop, \texttt{modpow}\\
$a/b < c/d$ ($b, d > 0$) & doubles & \texttt{a*d < c*b}\\
round a double & \texttt{(ll)x} & \texttt{llround(x)}\\
\bottomrule
\end{tabularx}
\code{math_int}

\topic{gcd, lcm and fractions}{$O(\log n)$}\label{t:gcd}
\use{Two repeating things line up every $\mathrm{lcm}(a, b)$. All $a_i$ equal mod $m$ iff
$m \mid \gcd(a_2 - a_1, a_3 - a_1, \dots)$. Steps of $k$ on a circle of $n$ visit
$n / \gcd(n, k)$ spots.}
\code{math_gcd}
\code{math_frac}

\topic{Extended Euclid: $ax + by = c$}{$O(\log n)$}
\use{Solvable iff $\gcd(a, b) \mid c$. $a x \equiv c \pmod m$ is $a x + m y = c$; with $c = 1$
it gives $1/a \bmod m$ for any $m$. Smallest $x \ge 0$ has the largest $y$: if that $y < 0$
there is no solution with $x, y \ge 0$.}
\code{math_egcd}

\topic{Modular arithmetic}{$O(\log e)$}\label{t:mod}
\use{\texttt{\% MOD} after every \texttt{+} and \texttt{*}; subtract as
\texttt{(a - b + MOD) \% MOD}; divide by multiplying with \texttt{inv(b)} (prime \texttt{MOD}).
Never compare values after \texttt{\%}. $a^e \equiv a^{e \bmod (p-1)} \pmod p$.}
\code{math_mod}

\subsection{Formulas and numbers}
\begin{tabularx}{\linewidth}{@{}l>{\raggedright\arraybackslash}X@{}}
\toprule
$1 + 2 + \dots + n$ & $n(n+1)/2$\\
$1^2 + 2^2 + \dots + n^2$ & $n(n+1)(2n+1)/6$\\
$1^3 + 2^3 + \dots + n^3$ & $\bigl(n(n+1)/2\bigr)^2$\\
$a + (a + d) + \dots$, $k$ terms & $k \cdot (\text{first} + \text{last}) / 2$\\
$1 + r + r^2 + \dots + r^{n-1}$ & $(r^n - 1)/(r - 1)$\\
$1 + \frac12 + \frac13 + \dots + \frac1n$ & $\approx \ln n + 0.58$\\
multiples of $k$ in $[L, R]$ & $\lfloor R/k \rfloor - \lfloor (L-1)/k \rfloor$\\
\midrule
primes $\le n$ & 25 ($10^2$), 168 ($10^3$), 78\,498 ($10^6$), 664\,579 ($10^7$)\\
most divisors & 240 ($10^6$), 1\,344 ($10^9$), 6\,720 ($10^{12}$), 103\,680 ($10^{18}$)\\
distinct primes & at most 9 ($n \le 10^9$), 15 ($n \le 10^{18}$)\\
prime moduli & $10^9 + 7$, $10^9 + 9$, $998\,244\,353$\\
\bottomrule
\end{tabularx}

\topic{Sieve of Eratosthenes}{$O(N \log \log N)$}\label{t:sieve}
\use{All primes up to $N$, or factorise many numbers $\le 10^7$ (0.1\,s, 40\,MB).}
\code{math_sieve}

\topic{Factorising one number}{$O(\sqrt n)$}
\use{$n \le 10^{12}$. From $n = \prod p_i^{e_i}$: divisors $\prod (e_i + 1)$, divisor sum
$\prod (p_i^{e_i + 1} - 1)/(p_i - 1)$, $\varphi(n) = n \prod (1 - 1/p_i)$, squarefree iff all
$e_i = 1$.}
\code{math_factor}

\topic{Divisors}{$O(\sqrt n)$}
\use{Try every divisor of $n$ (group sizes, periods). For all $x \le N$ at once: loop over the
multiples of every $d$, $O(N \log N)$.}
\code{math_divisors}

\topic{Chinese remainder theorem}{$O(\log m)$}\label{t:crt}
\use{$x \equiv a_1 \pmod{m_1}$, $x \equiv a_2 \pmod{m_2}$. Solutions repeat every
$\mathrm{lcm}(m_1, m_2)$; none if the moduli share a factor that the remainders disagree on.}
\code{math_crt}

\topic{Permutation cycles}{$O(n)$}\label{t:cycles}
\use{Repeated shuffle: after $k$ steps every element moves $k \bmod L$ along its cycle of
length $L$. Back to start after the lcm of the lengths. Sorting with swaps:
$n - \#\mathrm{cycles}$.}
\code{math_cycles}

\topic{Matrix power}{$O(k^3 \log n)$}\label{t:matrix}
\use{$f(n) = 2 f(n-1) + 3 f(n-2)$ for $n \le 10^{18}$: first row the coefficients, below it a
shifted identity; $f(n) = R[0][0] f(1) + R[0][1] f(0)$ for $R = M^{n-1}$. Walks of length $L$:
$A^L$.}
\code{math_matrix}

\topic{XOR tricks}{$O(n \log V)$}
\use{$x \oplus x = 0$; range xor via prefix xors; solve each bit on its own. Tree path
$u \to v$: $d[u] \oplus d[v]$. Graph walks: insert $d[a] \oplus d[b] \oplus w$ for every non-tree
edge into the basis.}
\code{math_xor}

\topic{Big numbers}{$O(\mathit{digits})$}\label{t:big}
\use{Hundreds of digits. Or submit that problem in Python 3: its integers never overflow.}
\code{math_bigint}
\pycode{bignum}

  \section{Counting and probability}

\subsection{Counting formulas}\label{t:counting}
\begin{tabularx}{\linewidth}{@{}Xl@{}}
\toprule
what you count & how many\\
\midrule
orders of $n$ different things & $n!$\\
pick $k$ of $n$, order matters & $n! / (n-k)!$\\
pick $k$ of $n$, order does not matter & $C(n, k)$\\
subsets of $n$ things & $2^n$\\
strings of length $k$ over $n$ letters & $n^k$\\
orders of a word, letter $i$ appears $c_i$ times & $n! / (c_1! \cdots c_m!)$\\
grid paths, $a$ right and $b$ up & $C(a+b, a)$\\
$x_1 + \dots + x_k = s$, all $x_i \ge 0$ & $C(s+k-1, k-1)$\\
$x_1 + \dots + x_k = s$, all $x_i \ge 1$ & $C(s-1, k-1)$\\
pick $k$ of $n$ kinds, repeats allowed & $C(n+k-1, k)$\\
$n$ people at a round table & $(n-1)!$\\
derangements $D_n$ (nobody in place) & $(n-1)(D_{n-1} + D_{n-2})$\\
Catalan: brackets, binary trees & $C(2n, n) / (n+1)$\\
\bottomrule
\end{tabularx}
\textbf{$D_n$:} 1, 0, 1, 2, 9, 44, 265. \textbf{Catalan:} 1, 1, 2, 5, 14, 42, 132, 429, 1430.

\topic{Binomial coefficients mod $p$}{$O(1)$ per query}
\use{$C(n, k)$ or $n!$ modulo a prime, $n \le 10^6$.}
\code{prob_binom}

\topic{Pascal's triangle}{$O(n^2)$}
\use{Small $n$, a modulus that is not prime, or $C(n, k)$ as a \texttt{double}.}
\code{prob_pascal}

\topic{Inclusion--exclusion}{$O(2^k \cdot k)$}
\use{At least one (or none) of $k \le 20$ properties: over every non-empty set $S$, $+$ the
count with all of $S$ if $|S|$ is odd, $-$ if even. Bounded stars and bars ($x_i \le u$):
$\sum_j (-1)^j C(k, j)\, C(s - j(u+1) + k - 1, k - 1)$.}
\code{prob_incl_excl}

\subsection{Probability rules}
\begin{tip}
\begin{itemize}
\item At least one $= 1 - P(\text{none})$, e.g.\ $1 - (1 - p)^n$.
\item Independent: $P(A \text{ and } B) = P(A) \cdot P(B)$.
  $P(A \mid B) = P(A \text{ and } B) / P(B)$.
\item Success with probability $p$ per try: on average $1/p$ tries.
\item \textbf{Linearity:} $E[X + Y] = E[X] + E[Y]$, always. Expected count
  $= \sum_{\text{items}} P(\text{item counts})$.
\item All $k$ sides of a die seen: $k(\frac1k + \frac1{k-1} + \dots + \frac11)$ throws.
\end{itemize}
\end{tip}

\topic{Expected-value DP}{states $\times$ moves}\label{t:expect}
\use{$E[s] = \mathit{cost}(s) + \sum_t P(s \to t) E[t]$, $E = 0$ at the end. Self-loop with
probability $q$: $E[s] = (c + \mathit{rest}) / (1 - q)$. You choose: take the $\min$.}
\code{prob_expect_dp}

\topic{Probability DP}{steps $\times$ states}
\use{Probability of every outcome after random steps: \texttt{nd[t] += dp[s] * q}. The
values must add up to 1.}
\code{prob_dp}

\topic{Answers as a fraction mod $p$}{}
\use{``Print $P \cdot Q^{-1} \bmod p$'': work mod $p$ throughout, $a/b$ becomes
\texttt{frac(a, b)}, divide by multiplying with \texttt{inv}.}
\code{prob_modfrac}

\subsection{Doubles}
\begin{tip}
\begin{itemize}
\item Print with \texttt{fixed <{}< setprecision(10)} (in the template).
\item Tiny probabilities ($0.5^{2000}$) underflow: add logarithms, \texttt{exp} at the end.
\item Unsure about a formula? Simulate $10^6$ times and compare on the samples.
\end{itemize}
\end{tip}

  \section{Games}

\subsection{Winning and losing}\label{t:games}
\begin{tip}
\begin{itemize}
\item Losing: every move leads to a winning position (or there is no move). Winning: some
  move leads to a losing one. Fill from the end positions.
\item Symmetric game: copying the other player's move often wins.
\item No idea? Print \texttt{win[]} for small $n$ and look for a pattern (often periodic).
\end{itemize}
\end{tip}

\topic{Win / lose table}{$O(n \cdot |S|)$}
\use{One pile, a move takes $s \in S$ stones, who cannot move loses.}
\code{game_win}

\topic{Nim and Grundy numbers}{$O(n \cdot |S|)$}
\use{Several independent games: xor of the Grundy numbers, 0 = you lose. Nim: xor of the
pile sizes; winning move: pile $p$ becomes $p \oplus x$ where $x$ is the xor.}
\code{game_grundy}

\topic{Taking from both ends}{$O(n^2)$}
\use{Players take the first or last element, each wants the most.}
\code{game_ends}

  \section{Geometry}

\subsection{Rules}\label{t:geo}
\begin{tip}
\begin{itemize}
\item Integer input: \texttt{T = ll}, stay exact (cross, squared lengths, \texttt{area2}).
  Coordinates up to $10^9$ fit. \texttt{sqrt} and division only at the end.
\item Doubles: never \texttt{==}; use \texttt{sgn}. \texttt{EPS = 1e-9} for coordinates up to
  $10^3$, larger coordinates need a larger \texttt{EPS}.
\item Distances: compare \texttt{norm(p - q) <= r * r}, not square roots.
\item Radians: \texttt{deg * M\_PI / 180}. Direction: \texttt{atan2(y, x)}, never
  \texttt{atan(y / x)}.
\item NaN: \texttt{acos(min(1.0, max(-1.0, v)))}, \texttt{sqrt(max(0.0, v))}.
\item Polygon: clockwise or not (sign of \texttt{area2})? First point repeated at the end?
\end{itemize}
\end{tip}

\topic{Point and cross product}{}
\use{Copy first. \texttt{dot} $> 0$: angle $< 90^\circ$, $= 0$: perpendicular.
\texttt{cross(a, b)} $> 0$: $b$ is left of $a$, $= 0$: parallel.}
\code{geo_point}

\topic{Segment intersection}{$O(1)$}
\use{Do two segments meet? Is a point on a segment? Exact with integers.}
\code{geo_segment}

\topic{Lines and distances}{$O(1)$}
\use{Crossing point of two lines, foot of a perpendicular, distance to a line or segment.}
\code{geo_line}

\topic{Polygon area, point in polygon}{$O(n)$}
\use{Shoelace area; inside / on border / outside. In triangle $abc$: \texttt{cross(a,b,q)},
\texttt{cross(b,c,q)}, \texttt{cross(c,a,q)} all $\ge 0$ or all $\le 0$.}
\code{geo_polygon}

\subsection{What shape is it?}
\begin{tip}
Use \emph{squared} lengths. \textbf{Triangle}, sides sorted $a \le b \le c$: right iff
$a + b = c$, acute iff $a + b > c$, obtuse iff $a + b < c$; collinear points are no triangle.
\textbf{Quadrilateral} $ABCD$, sides $s_1..s_4$ ($AB, BC, CD, DA$), diagonals $d_1 = AC$,
$d_2 = BD$. Most specific first:
\begin{itemize}
\item square: 4 sides equal, $d_1 = d_2$; rhombus: 4 sides equal
\item rectangle: $s_1 = s_3$, $s_2 = s_4$, $d_1 = d_2$; parallelogram: $s_1 = s_3$, $s_2 = s_4$
\item kite: ($s_1 = s_2$, $s_3 = s_4$) or ($s_2 = s_3$, $s_4 = s_1$)
\item trapezium: \texttt{cross(B-A, D-C) == 0} or \texttt{cross(C-B, A-D) == 0}
\end{itemize}
\end{tip}

\topic{Convex hull}{$O(n \log n)$}
\use{Smallest convex fence, its length or area. The two farthest points are hull vertices.}
\code{geo_hull}

\topic{Closest pair of points}{$O(n \log n)$}
\use{All pairs is fine up to $n \approx 10^4$. Random points: sort by $x$, look back only
while $(\Delta x)^2 <$ best.}
\code{geo_closest}

\topic{Angles and rotation}{}
\use{Direction \texttt{atan2(p.y, p.x)} in $(-\pi, \pi]$. Angle at $b$:
\texttt{angle(a - b, c - b)}. \texttt{byAngle}: exact sort by direction.}
\code{geo_angle}
\code{geo_rotate}

\topic{Circles}{}
\use{$p$ in circle: \texttt{norm(p - c) <= r * r}. Discs overlap iff
$|c_1 c_2|^2 \le (r_1 + r_2)^2$. Connected discs: test all pairs, then DSU.}
\code{geo_circle}

\subsection{Formulas}
\begin{tip}
\begin{itemize}
\item Pick: area $= I + B/2 - 1$, so $I = (|\texttt{area2}| - B + 2)/2$, with
  $B = \sum_{\text{edges}} \gcd(|\Delta x|, |\Delta y|)$.
\item $c^2 = a^2 + b^2 - 2ab\cos\gamma$, \quad $a / \sin\alpha = 2R$.
\item Heron: area $= \sqrt{s(s-a)(s-b)(s-c)}$, $s = (a+b+c)/2$. $R = abc / (4 A)$, $r = A / s$.
\item Polygon with $n$ corners: inside angles sum to $(n - 2) \cdot 180^\circ$.
\item Circle: $2\pi r$, $\pi r^2$, arc $r\theta$, sector $r^2\theta/2$, chord
  $2r\sin(\theta/2)$. Sphere: $4\pi r^2$, $\frac43\pi r^3$.
\end{itemize}
\end{tip}

  \section{Strings}

\subsection{Which string tool?}\label{t:strings}
\begin{tabularx}{\linewidth}{@{}Xll@{}}
\toprule
the problem & use & time\\
\midrule
anagrams, letter frequencies & count array & $O(n)$\\
letters form a palindrome? & parity mask & $O(n)$\\
longest palindromic substring & grow from centres & $O(n^2)$\\
find a pattern once & \texttt{s.find(p)} & $O(nm)$\\
all occurrences, period & KMP & $O(n + m)$\\
compare many substrings & hashing & $O(1)$ each\\
many words, prefix questions & trie & $O(\text{length})$\\
LCS, edit distance & DP chapter & $O(nm)$\\
\bottomrule
\end{tabularx}
\begin{tip}
$n \le 1000$: trying every substring with \texttt{substr} is often fine.
\end{tip}

\topic{Counting letters}{$O(n)$}
\use{Anagrams, ``can $A$ be made from the letters of $B$?''. Per position:
\texttt{cnt[i][c]}. In a range: prefix sums per letter.}
\code{str_count}

\topic{Palindromes}{$O(n^2)$}
\use{Longest palindrome inside a string: grow from each of the $2n - 1$ centres.}
\code{str_palindrome}

\topic{KMP}{$O(n + m)$}
\use{All occurrences of a pattern; the period of a string. $b$ is a rotation of $a$ iff
$|a| = |b|$ and $b$ occurs in $a + a$.}
\code{str_kmp}

\topic{Rolling hash}{$O(1)$ per substring}
\use{Compare many substrings. Longest repeated substring: binary search on the length, hashes
in a set.}
\code{str_hash}

\topic{Trie}{$O(\text{length})$}
\use{Many words, ``how many start with \dots?''. Without a trie: sort, the words starting with
$p$ are \texttt{lower\_bound(p)} to \texttt{lower\_bound(p + '\{')}.}
\code{str_trie}

\topic{Palindrome by rearranging}{$O(n)$}
\use{Possible iff at most one letter has an odd count: one bit per letter,
\texttt{mask \textasciicircum= 1 <{}< c}. Substring $s[l..r-1]$:
\texttt{pre[r] \textasciicircum{} pre[l]}.}
\code{str_parity}

  \section{Brute force and testing}

\topic{All subsets}{$O(2^n \cdot n)$}\label{t:subsets}
\use{$n \le 20$, choose some of the items. A few special things with 2 options each: try all
$2^k$.}
\code{bf_subsets}

\topic{All orders}{$O(n! \cdot n)$}\label{t:perms}
\use{$n \le 10$ and the answer is an order (route, schedule, assignment).}
\code{bf_perms}

\topic{Backtracking with pruning}{}
\use{Too many ways, but bad partial answers can be stopped early. Hardest items first; skip
options equal to one already tried.}
\code{bf_backtrack}

\topic{Stress test}{}
\use{WA and no idea why: \texttt{brute.cpp} (slow, surely right), \texttt{gen.cpp} (small
random tests), \texttt{sol.cpp}; run until the outputs differ. Crash or hang: the test is in
\texttt{in.txt}. Too slow? \texttt{time ./sol < big.txt} with a worst-case input.}
\code{bf_gen}
\shcode{stress}

\topic{Interactive problems}{}\label{t:interactive}
\use{Flush after every question. Test by typing the answers yourself, or with the testing
tool next to the samples.}
\code{bf_interactive}

\topic{Debugging}{}
\use{\texttt{dbg(x)} is in the template and only prints when compiled with \texttt{-DLOCAL} (the \texttt{c} command). \texttt{assert(cond)} stops on the line. \texttt{-D\_GLIBCXX\_DEBUG} checks every \texttt{v[i]}.}

\fi
\end{multicols*}
\ifdefined\onlysection\else
  % a ruled page for handwritten additions (allowed by the TRD rules)
  \clearpage
  \thispagestyle{notes}
  \section*{Notes}
  \addcontentsline{toc}{section}{Notes}
  {\color{black!25}\newcount\noteline
   \loop\ifnum\noteline<30 \vspace{5.4mm}\hrule\advance\noteline by 1 \repeat}
\fi
\end{document}
